La gestión del proyecto define cómo se organiza el equipo, en qué orden se realizan las actividades del Proyecto Fin de Curso y cómo se mantiene la comunicación con la organización. Las actividades se organizan a partir de las sesiones de elicitación planificadas y de las etapas de la ingeniería de requisitos descritas en la norma ISO/IEC/IEEE 29148 \cite{iso29148}.

\subsection{Roles del equipo}

\begin{table}[H]
  \centering\small
  \caption{Roles y responsabilidades del equipo en la entrega E1.}\label{tab:roles}
  \renewcommand{\arraystretch}{1.3}
  \begin{tabularx}{\textwidth}{@{}>{\raggedright\arraybackslash}p{3.6cm}>{\raggedright\arraybackslash}p{3.0cm}L@{}}
    \toprule
    Integrante & Rol & Responsabilidades en E1 \\
    \midrule
    Lavayen Bone Cristopher Javier & Integrador y analista líder & Estructura del repositorio y del documento, organización real, visión y alcance, manifiesto de evidencias, declaración de IA y compilación final. \\
    Aldas Zambrano Alan Andrés & Analista de contexto y ética & Diagrama de contexto, perspectivas de la organización, ética y tratamiento de datos, y modelo de consentimiento informado. \\
    Briones Solorzano William Adrian & Analista de partes interesadas & Identificación de partes interesadas, matriz poder--interés, estrategia de participación y registro de riesgos. \\
    Monserrate Ochoa Erick Ariel & Responsable de elicitación y planificación & Plan de elicitación, instrumentos (entrevista, cuestionario y guía de observación), observación del 18/11/2026 y gestión del proyecto. \\
    \bottomrule
  \end{tabularx}
\end{table}

\subsection{Cronograma hasta la semana 17}

El cronograma de la Tabla~\ref{tab:cronograma} toma como punto de partida las sesiones de la Tabla~\ref{tab:sesiones}. Las actividades desde la semana 7 son tentativas y se ajustarán a las fechas de entrega del sílabo y a la disponibilidad de la organización.

\begin{longtable}{@{}>{\raggedright\arraybackslash}p{1.3cm}>{\raggedright\arraybackslash}p{3.3cm}>{\raggedright\arraybackslash}p{6.8cm}>{\raggedright\arraybackslash}p{3.0cm}@{}}
  \caption{Cronograma del Proyecto Fin de Curso.}\label{tab:cronograma}\\
  \toprule
  Semana & Fechas & Actividad o entregable & Responsable \\
  \midrule
  \endfirsthead
  \toprule
  Semana & Fechas & Actividad o entregable & Responsable \\
  \midrule
  \endhead
  4  & 16--22 nov 2026 & Observación del turno (S-01) y entrega TA-PFC-E1 & Equipo \\
  5  & 23--29 nov 2026 & Entrevista al encargado del gimnasio (S-02) e inicio del cuestionario (S-03) & E.~Monserrate; C.~Lavayen, W.~Briones \\
  6  & 30 nov--6 dic 2026 & Cierre del cuestionario y reunión de validación con PI-01 (S-04) & A.~Aldas \\
  7  & 7--13 dic 2026 & Tabulación de resultados y consolidación de hallazgos & Equipo \\
  8  & 14--20 dic 2026 & Redacción de requisitos de usuario & Equipo \\
  9  & 21--27 dic 2026 & Clasificación y priorización de requisitos & Equipo \\
  10 & 28 dic 2026--3 ene 2027 & Revisión de avances y actualización del repositorio & Equipo \\
  11 & 4--10 ene 2027 & Modelado de procesos y casos de uso & Equipo \\
  12 & 11--17 ene 2027 & Especificación de requisitos (ERS) según ISO/IEC/IEEE 29148 \cite{iso29148} & Equipo \\
  13 & 18--24 ene 2027 & Validación de requisitos con PI-01 y PI-02 & Equipo \\
  14 & 25--31 ene 2027 & Correcciones a partir de la validación & Equipo \\
  15 & 1--7 feb 2027 & Prototipo de baja fidelidad y revisión con usuarios & Equipo \\
  16 & 8--14 feb 2027 & Integración del documento final del PFC & C.~Lavayen \\
  17 & 15--21 feb 2027 & Entrega y defensa del Proyecto Fin de Curso & Equipo \\
  \bottomrule
\end{longtable}

\subsection{Plan de comunicación}

\begin{table}[H]
  \centering\small
  \caption{Plan de comunicación del proyecto.}\label{tab:comunicacion}
  \renewcommand{\arraystretch}{1.3}
  \begin{tabularx}{\textwidth}{@{}LLlL@{}}
    \toprule
    Qué se comunica & A quién & Frecuencia & Medio \\
    \midrule
    Avance de tareas y acuerdos internos & Integrantes del equipo & Semanal & Reunión del equipo y grupo de mensajería \\
    Cambios en archivos del informe & Integrantes del equipo & En cada aporte & Commits en el repositorio de GitHub \\
    Avances del proyecto y solicitudes de validación & PI-01 Responsable del gimnasio en UBU & Por entrega & Correo institucional y reunión presencial \\
    Coordinación de entrevista y observaciones & PI-02 Encargado del gimnasio & Antes de cada sesión & Visita presencial al gimnasio \\
    Invitación a responder el cuestionario & PI-03, PI-04 y PI-05 Usuarios & Una vez (S-03) & Enlace al formulario institucional \\
    Avances y dudas del proyecto & Docente de la asignatura & Semanal & Clases y aula virtual (SGA) \\
    \bottomrule
  \end{tabularx}
\end{table}
